% ===================================================================
%  TAVOLA N. 8 — La mietitura. --- Caccia, vendemmia, pesca.
%  Traduction 2026 d'apres texto/io, controlee sur texto/fr et
%  texto/en. Consigne : voir texto/it/10-tavola-01.tex.
%
%  L'appel de note est dans le TITRE de la scene, non dans un alinea :
%  c'est ainsi que l'ido le compose.
%
%  POUR LA SECONDE FOIS EN QUATRE TABLEAUX, L'ACADEMIE IDO DELIBERE
%  SUR UN MOT QUE L'ITALIEN A DEPUIS SEPT CENTS ANS.
%
%  La note du titre dit ceci : « Parce que ce mot est imprecis : est-ce
%  une recolte de lin, une vendange, une cueillette de pommes ? Il
%  vaudrait peut-etre mieux employer MESONO, propose dans Mondo XIV,
%  p. 242. Mais jusqu'ici cette racine neuve n'a pas ete adoptee par
%  l'Academie et attend la discussion selon les regles habituelles du
%  mouvement ido. »
%
%  Or « mesono » est le latin « messis », la moisson, et « messis »
%  est l'italien messe, qui se lit dans Dante et dans tous
%  les dictionnaires. La racine n'est pas neuve : elle est vivante a
%  trois cents kilometres de Thaon-les-Vosges.
%
%  ET L'ITALIEN FAIT EXACTEMENT LA DISTINCTION QUE LA NOTE RECLAME.
%  Guignon se plaint qu'un seul mot serve pour trois choses ;
%  l'italien en a trois :
%     mietitura  la moisson des cereales, et rien d'autre ;
%     vendemmia  la recolte du raisin, et rien d'autre ;
%     raccolta   le terme general, quand on veut le general.
%  Les deux scenes de cette planche sont justement l'une une
%  mietitura et l'autre une vendemmia, et la colonne italienne est la
%  seule des trente-quatre a pouvoir les nommer sans hesiter.
%
%  C'est le meme constat qu'au tableau 5, ou la note sur « balk-o »
%  citait « travicello » sans voir qu'il repondait a la question. DEUX
%  FOIS L'ACADEMIE CHERCHE UN MOT QUI EST DANS SA PROPRE LISTE. La
%  raison est peut-etre simple : ses membres lisaient l'italien dans
%  les dictionnaires et non dans les champs.
%
%  ET « VENDEMMIA » DIT ENCORE AUTRE CHOSE. Le mot est « vini-demia »,
%  l'enlevement du VIN. On ne recolte pas le raisin : on recolte le
%  vin, deux mois avant qu'il existe. La chose est nommee par ce
%  qu'elle deviendra, et c'est le contraire du mosto de
%  l'alinea 3, qui est nomme par ce qu'il est encore.
%
%  Le tableau 8 allemand avait trouve que tout le vocabulaire allemand
%  de la vigne etait latin, et que la frontiere du vin et de la biere
%  etait celle de l'empire. Cette page-ci est de l'autre cote de cette
%  frontiere : elle n'emprunte rien parce qu'elle est le pret.
% ===================================================================

%%K t08-noto-1 noto
\VUnotes{143.2mm}{%
\textit{(*) Perché questa parola è imprecisa: è una raccolta di lino, una vendemmia,
una raccolta di mele? Forse sarebbe meglio usare «}
\VUgras{mesono} \textit{», proposto in} Mondo \textit{XIV, p. 242. Finora però
questa radice nuova non è stata adottata dall'Accademia e attende la discussione
secondo le regole consuete del movimento ido.}%
}

%%K t08-apar-1 apar
\VUsaut{5.0mm}
\VUtitre{15.0pt}{60}{TAVOLA N. 8}
\VUsaut{3.0mm}
\VUcentre{13.2pt}{90}{\textit{Prima scena.}}
\VUsaut{3.0mm}
\VUpk{10.2pt}{40}{La mietitura (*).}
\VUsaut{4.0mm}
\VUpk{10.2pt}{40}{Il volto della campagna.}
\VUsaut{3.0mm}

%%K t08-c1-01-1 p
1. --- Con l'\VUgras{estate} la campagna ha cambiato di nuovo volto. La semente
affidata al solco è germogliata; una piccola pianta verde è uscita dalla terra e ha
portato la sua spiga. Il chicco si è formato, poi è maturato, e adesso non resta che
raccogliere il grano.

%%K t08-c1-02-1 p
2. --- I \VUgras{fiori di campo} sono sbocciati. \VUgras{Fiordalisi}
\textsuperscript{(1)}, \VUgras{papaveri} \textsuperscript{(2)} e
\VUgras{digitali} \textsuperscript{(3)} spiccano sul tono uniforme dei campi di grano
con il contrasto allegro delle loro \VUgras{corolle} fresche.

%%K t08-c1-03-1 p
3. --- Gli alberi da frutto hanno messo le foglie; la frutta è matura. Il piccolo
\VUgras{ciliegio} \textsuperscript{(4)} è carico di belle \VUgras{ciliegie}
\textsuperscript{(5)}, che i bambini colgono e mangiano con gusto.

%%K t08-c1-04-1 p
4. --- Il gregge può ormai restare fuori tutto il giorno.

%%K t08-c1-04-2 p
Gli \VUgras{agnelli} \textsuperscript{(6)} sono cresciuti e sono diventati belle
\VUgras{pecore} \textsuperscript{(7)} e bei \VUgras{montoni} \textsuperscript{(8)} dal
vello folto. Pascolano tranquilli l'\VUgras{erba} del \VUgras{pascolo}
\textsuperscript{(9)}, punteggiato di \VUgras{margherite}.

%%K t08-c1-04-3 p
Presto questi animali saranno tosati per la loro \VUgras{lana}, con cui si fa il panno
dei nostri vestiti.

%%K t08-tit-2 sub
\VUsaut{5.0mm}
\VUpk{10.2pt}{40}{Il mandriano.}
\VUsaut{3.4mm}

%%K t08-c1-05-1 p
5. --- Il \VUgras{mandriano} \textsuperscript{(10)} non sembra badare molto a loro. La
sua unica preoccupazione è il temporale che passa all'orizzonte, e il \VUgras{pastore}
\textsuperscript{(13)}, che alza la sua \VUgras{borraccia} \textsuperscript{(14)} alle
labbra, sembra ancora più spensierato.

%%K t08-c1-06-1 p
6. --- Il caldo è soffocante: perfino il \VUgras{cane} \textsuperscript{(15)} viene a
rinfrescarsi; lecca l'acqua fresca del \VUgras{ruscello} \textsuperscript{(16)} e fa
scappare una povera \VUgras{rana} \textsuperscript{(17)} che stava accovacciata
nell'erba della riva. Salta nell'acqua limpida dove fioriscono le
\VUgras{ninfee} \textsuperscript{(18)}.

%%K t08-c1-07-1 p
7. --- Una \VUgras{cutrettola} \textsuperscript{(19)} fa il bagno accanto alla piccola
\VUgras{sorgente} \textsuperscript{(20)} che sgorga fra due rocce e si nasconde sotto
gli \VUgras{sterpi spinosi} \textsuperscript{(21)} e i \VUgras{biancospini}
\textsuperscript{(22)}.

%%K t08-tit-3 sub
\VUsaut{5.0mm}
\VUpk{10.2pt}{40}{La mietitura.}
\VUsaut{3.4mm}

%%K t08-c1-08-1 p
8. --- Per i bambini di campagna la mietitura del grano è un periodo molto divertente.
Fanno allegre \VUgras{capriole} \textsuperscript{(24)} sui \VUgras{mucchi di paglia}
\textsuperscript{(23)}, aiutano i \VUgras{mietitori} \textsuperscript{(26)}, e mentre
questi tagliano il \VUgras{campo di grano} \textsuperscript{(27)} con le loro
\VUgras{falci} \textsuperscript{(25)}, ben affilate sulla \VUgras{cote}
\textsuperscript{(30)}, loro raccolgono il grano e lo legano in \VUgras{covoni}
\textsuperscript{(31)} o in piccoli \VUgras{fasci} \textsuperscript{(32)}.

%%K t08-c1-09-1 p
9. --- A volte ai più grandi si permette di usare una \VUgras{falce messoria}
\textsuperscript{(12)} sugli \VUgras{steli dorati} \textsuperscript{(28)} che portano
le pesanti \VUgras{spighe} \textsuperscript{(29)} piene di chicchi; e i piccoli,
mentre sorvegliano il \VUgras{bestiame} \textsuperscript{(33)} che pascola nel
\VUgras{prato} \textsuperscript{(34)} all'ombra del grande \VUgras{tiglio}
\textsuperscript{(35)}, prendono le belle \VUgras{farfalle} con la loro \VUgras{rete}
\textsuperscript{(36)}.

%%K t08-c1-09-2 p
Qualcuno sale perfino sul \VUgras{carro} \textsuperscript{(37)} per sistemare i covoni
che gli vengono passati su con un \VUgras{forcone} \textsuperscript{(38)}.

%%K t08-c1-10-1 p
10. --- Su tutta la \VUgras{pianura} \textsuperscript{(39)}, dove si alzano in volo le
\VUgras{allodole} \textsuperscript{(41)}, c'è un gran movimento. Sono tutti occupati
con la mietitura. Ma mentre i poveri continuano a usare gli attrezzi di una volta ---
\VUgras{falci}, \VUgras{falci messorie} e \VUgras{correggiati}
\textsuperscript{(40)} --- i proprietari ricchi fanno tagliare il loro frumento o la
loro \VUgras{segale} \textsuperscript{(43)} con una \VUgras{mietitrice}
\textsuperscript{(42)}. Poi, senza dover portare i covoni nel fienile, si fanno grandi
\VUgras{biche} \textsuperscript{(44)} sul posto.

%%K t08-c1-11-1 p
11. --- Poi si porta una \VUgras{trebbiatrice} \textsuperscript{(47)}, azionata da una
\VUgras{locomobile} \textsuperscript{(45)} tramite una \VUgras{cinghia}
\textsuperscript{(46)}, e così il grano viene separato dalla \VUgras{paglia} senza
lasciare il campo in cui è stato seminato.

%%K t08-tit-4 sub
\VUsaut{5.0mm}
\VUpk{10.2pt}{40}{Il temporale.}
\VUsaut{3.4mm}

%%K t08-c1-12-1 p
12. --- Al tempo della mietitura scoppiano spesso violenti \VUgras{temporali}
\textsuperscript{(53)}. Prima si sente da lontano il brontolio del \VUgras{tuono}; un
forte \VUgras{vento di ponente} \textsuperscript{(54)} si alza e piega gli alberi più
grandi del \VUgras{bosco} \textsuperscript{(49)} finché non gemono. Nere
\VUgras{nuvole} \textsuperscript{(56)} invadono il cielo; sono attraversate dagli
zigzag accecanti del \VUgras{lampo} \textsuperscript{(55)}. La \VUgras{pioggia}, e a
volte la \VUgras{grandine}, comincia a cadere e schiaccia il raccolto che sta per
essere tagliato.

%%K t08-c1-13-1 p
13. --- Spesso il fulmine dà fuoco a un edificio. L'anno scorso, per esempio, il tetto
del \VUgras{mulino a vento} \textsuperscript{(50)} è bruciato del tutto e due delle sue
\VUgras{ali} \textsuperscript{(51)} sono andate in pezzi.

%%K t08-c1-14-1 p
14. --- Dopo qualche ora torna la calma. Un \VUgras{arcobaleno} \textsuperscript{(52)}
splendente appare all'orizzonte, e la natura riprende la vita che è stata interrotta
per un momento. I contadini, che si erano riparati nelle \VUgras{capanne}
\textsuperscript{(48)}, tornano al lavoro e si mettono di buona lena a riparare i
danni che hanno appena subito.

%%K t08-tit-5 sub
\VUsaut{6.0mm}
\VUcentre{13.2pt}{90}{\textit{Seconda scena.}}
\VUsaut{3.6mm}

%%K t08-tit-6 sub
\VUpk{10.2pt}{40}{Caccia. --- Vendemmia.}
\VUsaut{1.4mm}
\VUpk{10.2pt}{40}{Pesca.}
\VUsaut{4.0mm}

%%K t08-c2-01-1 p
1. --- Verso la fine di \VUgras{agosto} o a metà \VUgras{settembre}, secondo la
regione, si apre la caccia. Appena i primi raggi di sole hanno fatto uscire le
\VUgras{lucertole} \textsuperscript{(2)} dai loro buchi, il \VUgras{cacciatore}
\textsuperscript{(5)} parte con i suoi \VUgras{cani} \textsuperscript{(4)}.

%%K t08-c2-02-1 p
2. --- Si è messo il suo robusto \VUgras{abito da caccia} \textsuperscript{(9)} di
panno pesante e le \VUgras{ghette} di cuoio che gli permettono di camminare nell'erba
bagnata dove si nascondono i \VUgras{rospi} \textsuperscript{(1)}; ha preso il suo
\VUgras{fucile} \textsuperscript{(6)} e il suo \VUgras{carniere}
\textsuperscript{(10)}, ed eccolo che parte.

%%K t08-c2-03-1 p
3. --- La foga della caccia lo porta in mezzo a una vigna dove la vendemmia è in pieno
svolgimento. I figli del vignaiolo, che di solito guardano le capre sul ciglio della
strada, oggi le lasciano pascolare dove vogliono, e dopo aver legato a un palo un
vecchio \VUgras{caprone} \textsuperscript{(3)} che approfitterebbe certamente della
libertà per scappare, si prendono anche loro la loro parte di divertimento. Uno prende
un grappolo da un \VUgras{cesto della vendemmia} \textsuperscript{(12)} e si diverte a
spremere il \VUgras{succo} \textsuperscript{(14)} in una \VUgras{tazza}
\textsuperscript{(13)} per fare il mosto (vino non fermentato). Un altro si sta
mettendo in testa una \VUgras{corona di pampini} \textsuperscript{(15)} che ha appena
intrecciato. Vicino a loro, \VUgras{passeri} \textsuperscript{(16)} sfacciati arrivano
fino al torchio per rubare l'uva.

%%K t08-c2-04-1 p
4. --- Mentre i bambini si divertono, gli uomini lavorano. I \VUgras{vendemmiatori}
\textsuperscript{(18)} portano l'uva in cantina in \VUgras{gerle}
\textsuperscript{(17)}; un \VUgras{bottaio} \textsuperscript{(19)} ripara le
\VUgras{botti} \textsuperscript{(20)}, controlla che le \VUgras{doghe}
\textsuperscript{(22)} e i \VUgras{fondi} \textsuperscript{(21)} siano sani, e
sostituisce i \VUgras{cerchi} \textsuperscript{(23)} rotti. È stato lui a fare anche i
grandi \VUgras{tini} \textsuperscript{(25)} in cui si versa l'\VUgras{uva}
\textsuperscript{(26)} a fermentare.

%%K t08-c2-05-1 p
5. --- Quando la fermentazione è finita, il vino viene tolto e il resto si mette al
\VUgras{torchio} \textsuperscript{(28)}; stringendo poco a poco la grande
\VUgras{vite} \textsuperscript{(29)}, si spreme il succo, che cola in un
\VUgras{mastello} \textsuperscript{(32)} e dà ancora altro \VUgras{vino}
\textsuperscript{(31)}.

%%K t08-tit-8 sub
\VUsaut{6.0mm}
\VUpk{10.2pt}{40}{Birra e sidro.}
\VUsaut{3.4mm}

%%K t08-c2-06-1 p
6. --- Sul muro della \VUgras{cantina} \textsuperscript{(27)} si arrampica un ramo di
\VUgras{luppolo} \textsuperscript{(30)}. Qui è solo una pianta ornamentale, ma in certi
paesi, dove la vite è molto rara, si coltiva il luppolo per fare la \VUgras{birra}.

%%K t08-c2-07-1 p
7. --- Il piccolo \VUgras{melo} \textsuperscript{(33)} è carico di mele con cui si farà
un ottimo \VUgras{sidro} quando saranno mature. Nella loro \VUgras{gabbia}
\textsuperscript{(34)} il \VUgras{canarino} \textsuperscript{(35)} e il
\VUgras{cardellino} \textsuperscript{(36)} guardano con invidia i passeri che
saltellano in libertà.

%%K t08-c2-08-1 p
8. --- A perdita d'occhio, sul pendio delle colline, ci sono file di \VUgras{pali della
vigna} \textsuperscript{(40)} che sostengono i magnifici \VUgras{ceppi di vite}
\textsuperscript{(39)}, carichi di \VUgras{grappoli}.

%%K t08-c2-08-2 p
Tutto l'anno il \VUgras{vignaiolo} \textsuperscript{(42)} ha lavorato per curare la
sua \VUgras{vigna} \textsuperscript{(38)}, e adesso è ricompensato della sua fatica con
un buon raccolto. Le \VUgras{vendemmiatrici} \textsuperscript{(41)} riempiono i tini
posati sul carro tirato dai \VUgras{muli} \textsuperscript{(43)}, e sono tutti allegri.

%%K t08-tit-9 sub
\VUsaut{5.0mm}
\VUpk{10.2pt}{40}{Pesca.}
\VUsaut{3.4mm}

%%K t08-c2-09-1 p
9. --- Lo stesso vale per i \VUgras{pescatori con la lenza} \textsuperscript{(44)}, che
sono venuti alle prime luci a sistemarsi in riva all'acqua con le loro
\VUgras{canne da pesca} \textsuperscript{(45)}.

%%K t08-c2-09-2 p
I \VUgras{pesci} \textsuperscript{(47)} abboccano appena calano la loro \VUgras{lenza}
\textsuperscript{(46)} nell'acqua, e non hanno quasi il tempo di rimettere l'\VUgras{esca}
che una nuova vittima si dibatte in fondo alla lenza.

%%K t08-c2-09-3 p
Però il \VUgras{pescatore con la rete} \textsuperscript{(48)}, che getta la sua rete
dalla sua \VUgras{barca} \textsuperscript{(49)} in mezzo al \VUgras{fiume}
\textsuperscript{(50)}, ne prende molti di più.

%%K t08-c2-10-1 p
10. --- L'autunno è la stagione delle belle gite: a piedi, a \VUgras{cavallo}
\textsuperscript{(52)}, in \VUgras{bicicletta} \textsuperscript{(53)}, in
\VUgras{automobile} \textsuperscript{(54)}. Le \VUgras{strade} \textsuperscript{(51)}
sono pulite e si approfitta delle ultime belle giornate, perché le \VUgras{rondini}
\textsuperscript{(56)} si radunano già sui tetti per partire verso paesi più caldi. Le
foglie del vecchio \VUgras{noce} \textsuperscript{(57)} ingialliscono, le
\VUgras{noci} cadono, e gli alti \VUgras{alberi} riuniti in un \VUgras{gruppo}
\textsuperscript{(58)} sull'\VUgras{altura} \textsuperscript{(59)} saranno presto
spogli.

%%K t08-c2-11-1 p
11. --- Il \VUgras{sole} \textsuperscript{(60)} sorge più tardi, le giornate si
accorciano, al \VUgras{mattino} si sente già un fresco abbastanza pungente, e voli di
\VUgras{beccacce} \textsuperscript{(61)} lasciano già il nostro clima inospitale.
\VUsaut{6.0mm}
\VUfilet{22mm}
